\section{Introduction}
\label{sec:intro}

\IEEEPARstart{I}{nspecting} a harbor, a dam, or an offshore platform requires information both
above and below the water surface. Aerial robots survey the superstructure, ground robots
cover the quay, surface vessels patrol the waterline, and underwater vehicles inspect piles,
hulls, and foundations. The value of such a heterogeneous team lies in a \emph{single},
consistent map. Yet collaborative SLAM (C-SLAM) research has progressed \emph{within} a medium:
decentralized air/ground systems share maps over radio~\cite{tian2022kimeramulti,
lajoie2024swarmslam, liu2025slideslam}, and underwater teams share compact sonar maps over
acoustic modems~\cite{huang2025draco2}. The first cross-medium C-SLAM system links a surface
vessel and underwater vehicles through structures visible above and below the water, using a
centralized back-end~\cite{mcconnell2026aboveandbelow}.

\draftnote{Paragraph: why the waterline is hard. (1) Perception split: the same pile is seen
by LiDAR/cameras above and by imaging sonar below; the two observations share only the
horizontal position (Fig.~\ref{fig:coaxial}). (2) Communication split: RF links carry
megabits per second but do not propagate in water; acoustic links carry $10^2$--$10^4$~bit/s
with seconds of latency and frequent loss; an AUV reaches RF only when surfaced. (3) No
central server can be assumed reachable by all agents at all times.}

\draftnote{Paragraph: consequences for design: decentralization, information that is cheap to
send, double-counting-free sharing, cross-medium association that is robust to aliasing in
regular structures (pile rows).}

\noindent This paper makes the following contributions:
\begin{enumerate}
  \item \avatar{}, the first decentralized C-SLAM system for teams spanning aerial, ground,
        surface, and underwater agents. Each agent keeps a 4-DoF factor graph on a shared
        waterline datum and shares condensed landmarks computed from its own measurements
        only, which rules out double counting by construction (Sec.~\ref{sec:system}).
        \todo{confirm novelty ledger N1.}
  \item A cross-medium association method: coaxial landmark parts, uncertainty-scaled
        gating, and a pairwise-consistency graph with an explicit ambiguity test. The method
        yields direct UAV--AUV constraints through waterline-crossing structures
        (Sec.~\ref{sec:association}). \todo{ledger N3 after reading~\cite{mcconnell2026aboveandbelow} in full.}
  \item A medium-aware communication layer: link models, a 16-byte-per-landmark wire format,
        and \todo{VoI-per-byte scheduler} across RF and acoustic links (Sec.~\ref{sec:comm}).
  \item An open benchmark: a fast multi-domain simulator and ROS~2/Gazebo scenarios with
        ground truth and communication traces (Sec.~\ref{sec:experiments}).
\end{enumerate}

\begin{figure}[t]
  \centering
  \fbox{\parbox{0.92\columnwidth}{\centering\vspace{2em}\todo{Teaser: harbor with UAV, UGV,
  USV, AUVs; above/below parts of a pile; RF vs.\ acoustic links.}\vspace{2em}}}
  \caption{\avatar{} overview. \todo{caption}}
  \label{fig:teaser}
\end{figure}
